% Source: https://artofproblemsolving.com/community/c3h2009671_infinite_series
% Stable identity: aops-topic-2009671-q01
Adding exponents,
\[P_N=i^{1+2+\cdots+N}=i^{T_N},\qquad T_N=\dfrac{N(N+1)}{2}.\]
Since $i^4=1$, $P_N$ depends only on $T_N\bmod 4$. Since $T_{N+8}-T_N=8N+36\equiv0\pmod4$, the value of $T_N\bmod4$, and hence $P_N$, is periodic in $N$ with period $8$. For $N=1,\dots,8$:
\[T_N=1,3,6,10,15,21,28,36\equiv1,3,2,2,3,1,0,0\pmod4,\]
\[P_N=i,\ -i,\ -1,\ -1,\ -i,\ i,\ 1,\ 1,\]
and the pattern repeats every $8$ terms.

So $P_N$ takes each of the values $\boxed{1,\ i,\ -1,\ -i}$ (each twice in every period of $8$). The partial products keep jumping among four different numbers and never settle down, so they have no limit: the infinite product $\displaystyle\prod_{n=1}^{\infty}i^n$ \textbf{does not converge}. (Every term has absolute value $1$ and does not approach $1$, which already rules out convergence.)
